\scanpage{86}
\mspage{92}

donc ($n$ étant inversible dans $k$) :

\[
H^{\unsure{\cdot|i}}(X,k)=R\Gamma_X(k_X)\simeq\deleted{\unsure{R}}H^{\cdot}(Z,k)^G\unsure{\relbar\joinrel\relbar|\longrightarrow}H^*(Z,k)\ ,
\]

ce n'est en général pas un iso. On peut prendre p. ex. $G=\mathbb Z/2\mathbb Z$, $Z\sim S^1$ comme avant, et $G$ opérant par symétrie de $S^1$, comme précédemment.

Cet exemple est particulièrement convainquant, car c'est un exemple \uline{avec} $S,S'$ $\mathbb D$-\uline{asphériques}, et $S'\to S$ un iso. local (en fait, un revêtement principal de $S$).

C'est aussi un exemple d'un $\mathbb D$-fibré lisse \unsure{et|plus} propre $X\to S$, sur une base \uline{asphérique} $S$, et tel que l'inclusion d'une fibre dans $X$

\[
X_s\unsure{\relbar|\longrightarrow}X
\]

ne soit \uline{pas} une $\mathbb D$-équivalence.

Le seul reproche qu'on puisse faire à cet exemple, en comparaison avec le précédent, c'est que les catégories $X,S$ ne sont pas \inserted{des ens.} ordonnés. \inserted{Et} $S$ ne peut être un ensemble ordonné fini, car l'espace classifiant d'un groupe fini a de la cohomologie en toutes dimensions (sauf erreur).

% archive sheet number: 85
